%% kapa: exam=1 round=30 year=2019 period=2 form=A subject=accounting subject_ko=회계학 seq=19 number=59
%% source: Q-Net 감정평가사 기출문제 (cst00309) artlSeq=5205205 file=extracted/1008/5205205/2019년 제30회 감정평가사 1차시험 문제지/2019년 제30회 감정평가사 1차시험 문제지/2019년도 제30회 감정평가사 1차 2교시 A형.hwp
%% conversion: kordoc 4.24.51 -> tools/gen1.py (auto); review=proofread-sonnet; flags=none
\begin{question}{59}
재무보고를 위한 개념체계에 관한 설명으로 옳은 것은?
\begin{choices}
  \item 부채는 과거 거래의 결과로 발생한 것으로 미래에 기업실체가 부담할 의무로서 현재 그 의무를 이행한다.
  \item 미래경제적효익의 유입 가능성이 높거나 그 항목의 원가를 신뢰성 있게 측정할 수 있으면 재무제표에 자산으로 인식한다.
  \item 비용은 경제적 효익이 유출, 소비됨으로써 자산이 증가하거나 부채가 감소할 가능성이 높고 그 금액을 신뢰성 있게 측정할 수 있을 때 인식한다.
  \item 자산의 사용정도를 체계적이고 합리적인 방법으로 배분하는 감가상각비는 관련 수익과의 관련성이 직접적으로 파악, 결정되는 비용이다.
  \item 표현충실성을 위한 서술에 오류가 없다는 것은 현상의 기술에 오류나 누락이 없고, 보고 정보를 생산하는데 사용되는 절차의 선택과 적용에 절차상의 오류가 없음을 의미한다.
\end{choices}
\end{question}
\answer{5}
